\documentclass[10pt,a4paper]{amsart}
\usepackage[margin=19mm]{geometry}
\usepackage{amsmath,amssymb,mathtools}
\usepackage[scr=rsfs,cal=euler]{mathalfa}
\usepackage{xfrac}
\usepackage[unicode,hidelinks,bookmarks=false]{hyperref}
\newcommand{\ie}{i.e.\ }
\newcommand{\cf}{cf.\ }
\newcommand{\resp}{respectively\ }
\newcommand{\Subsection}{\subsection}
\providecommand{\nobreakdash}{\nobreak\hbox{-}}
\setlength{\emergencystretch}{2em}
\multlinegap0pt
\def\C{{\mathbb{C}}}
\def\MM{{\mathbf{M}}}
\def\bi{\hbox{\bf{i}}}
\def\cB{{\mathcal{B}}}
\def\sC{{\mathscr{C}}}
\def\sM{{\mathscr{M}}}
\def\sP{{\mathscr{P}}}
\theoremstyle{plain}
\newtheorem{theorem}{Theorem}
\newtheorem{proposition}[theorem]{Proposition}
\newtheorem{lemma}[theorem]{Lemma}
\newtheorem{corollary}[theorem]{Corollary}
\theoremstyle{definition}
\newtheorem{definition}[theorem]{Definition}
\newtheorem{assumption}[theorem]{Assumption}
\theoremstyle{remark}
\newtheorem{remark}[theorem]{Remark}
\title[Scalable Ground-State Certification of Quantum Spin Systems via Structured Noncommutative Polynomial Optimization]{Scalable Ground-State Certification of Quantum Spin Systems via Structured Noncommutative Polynomial Optimization: Proposition 4.1}
\author{Jie Wang, David Jansen, Ir\'en\'ee Fr\'erot, Marc-Olivier Renou, Victor Magron, Antonio Ac\'{\i}n}
\date{Source: arXiv:2604.01555v2 (CC BY 4.0)}
\begin{document}
\maketitle

\noindent\emph{Source and licence.} The delivered work is the article shown in the title above, version arXiv:2604.01555v2 (CC BY 4.0), distributed under the Creative Commons Attribution 4.0 International licence, as recorded in the arXiv licence metadata for this exact version and shown on the abstract page saved with this item. The full licence notice, which carries the licence URL and the address of that abstract page, is the bundled file \texttt{licenses/arxiv-2604.01555-CC-BY-4.0.txt}.\par\smallskip

\noindent\textit{Editorial scope.} Counted unit: the article's Proposition 4.1 (source0.tex lines 655--657, source label \texttt{prop2}); statement and proof are the article's own text line for line, in the article's own notation ($\MM_d(\ell)$, $\cB_d^{(i,j)}$, $\eta(u)$, $\ell(u)$), with no line omitted, substituted or re-flowed. Required context retained uncounted: the article's four subbasis definitions and the moment-relaxation display (source0.tex lines 301--312 and 386--394), its Lemma on the real part of the relaxation (lines 598--605) and its block-matrix lemma (lines 615--627). Editorial formatting only: an AMS \texttt{amsart} preamble that restates the article's own macro definitions (\texttt{\textbackslash C}, \texttt{\textbackslash MM}, \texttt{\textbackslash bi}, \texttt{\textbackslash cB}, \texttt{\textbackslash sC}, \texttt{\textbackslash sM}, \texttt{\textbackslash sP}) and the theorem environments the retained lines use; the article's SIAM class file \texttt{siamonline220329.cls} is neither shipped nor input; sequential numbering of the retained environments in order of appearance, whereas the article prints them with its own section-relative numbers. No citation occurs inside any retained line, so the wrapper carries no bibliography. No asset-inclusion command, no drawing environment and no image asset occurs in this package.


\section*{Retained context: rules (source0.tex lines 301--312)}

\begin{subequations}\label{rules}
\begin{align}
    (\sigma_i^a)^2\quad&\longrightarrow\quad1,\quad i=1,\ldots,N,a\in\{x,y,z\},\\
    \sigma_i^x\sigma_i^y\quad&\longrightarrow\quad\bi\sigma_i^z,\quad i=1,\ldots,N,\label{rule1}\\
    \sigma_i^y\sigma_i^x\quad&\longrightarrow\quad-\bi\sigma_i^z,\quad i=1,\ldots,N,\\
    \sigma_i^y\sigma_i^z\quad&\longrightarrow\quad\bi\sigma_i^x,\quad i=1,\ldots,N,\\
    \sigma_i^z\sigma_i^y\quad&\longrightarrow\quad-\bi\sigma_i^x,\quad i=1,\ldots,N,\\
    \sigma_i^z\sigma_i^x\quad&\longrightarrow\quad\bi\sigma_i^y,\quad i=1,\ldots,N,\\
    \sigma_i^x\sigma_i^z\quad&\longrightarrow\quad-\bi\sigma_i^y,\quad i=1,\ldots,N,\label{rule6}\\
    \sigma_i^a\sigma_j^b\quad&\longrightarrow\quad\sigma_j^b\sigma_i^a,\quad 1\le i\ne j\le N,a,b\in\{x,y,z\}.
\end{align}
\end{subequations}


\section*{Retained context: mom (source0.tex lines 386--394)}

\begin{equation}\label{mom}
\begin{aligned}
\lambda_d^*\coloneqq\min\limits_{\{\ell(u)\}_{u\in\tilde{W}_{2d}}}&\quad\ell(H)\\
    \text{subject to:}&\quad\MM_{d}(\ell)\succeq0, \\
    &\quad\MM_{d}(\ell) \text{ obeys the monomial replacement rules \eqref{rules}},\\
    &\quad\ell(u^{\star})=\overline{\ell(u)},\quad \forall u\in\tilde{W}_{2d},\\
    &\quad\ell(1)=1.
\end{aligned}
\end{equation}


\section*{Retained context: lm1 (source0.tex lines 598--605)}

\begin{lemma}\label{lm1}
Let $A,B\in\C^{n\times n}$ such that $A^*=A$ and $B^*=-B$. Then
\begin{equation}
A+B\bi\succeq0\iff\begin{bmatrix}
    A&B\\-B&A
\end{bmatrix}\succeq0.
\end{equation}
\end{lemma}


\section*{Retained context: lm2 (source0.tex lines 615--627)}

\begin{lemma}\label{lm2}
Let $A\in\C^{n\times n}, B\in\C^{m\times m}, C\in\C^{m\times n}$ such that $A^*=A$ and $B^*=B$. Then,
\begin{equation}
\begin{bmatrix}
A&C\bi\\
-C^{*}\bi&B
\end{bmatrix}\succeq0\iff
\begin{bmatrix}
A&C\\
C^{*}&B
\end{bmatrix}\succeq0.
\end{equation}
\end{lemma}


\section{The counted unit: Proposition 4.1 and its complete original proof}

\begin{proposition}\label{prop2}
In the moment relaxation \eqref{mom} for Heisenberg models, there is no loss of generality in assuming that $\ell(u)=0$ whenever $\eta(u)\ne(1,1,1)$.
\end{proposition}

\begin{proof}
For $i\in\{1,2,3,4\}$, let $\MM_{d}^{(i)}(\ell)$ be the block of the moment matrix $\MM_{d}(\ell)$ indexed by the subbasis $\cB_d^{(i)}$. For each $i$, we define the matrices $A^{(i)},B^{(i)}\in\C^{|\cB_d^{(i,1)}|\times|\cB_d^{(i,1)}|},C^{(i)},D^{(i)}\in\C^{|\cB_d^{(i,1)}|\times|\cB_d^{(i,2)}|},E^{(i)},F^{(i)}\in\C^{|\cB_d^{(i,2)}|\times|\cB_d^{(i,2)}|}$ as follows:
\begin{subequations}\label{sec:eq2}
\begin{align}
    A^{(i)}_{v,w}&=\begin{cases}
    \sC(v^{\star}w)\ell(\sM(v^{\star}w)), \quad&\text{if }2\mid\deg(\sM(v^{\star}w)),\\
    0, \quad&\text{otherwise},
    \end{cases}\quad \text{where } v,w\in\cB_d^{(i,1)};\\
    B^{(i)}_{v,w}&=\begin{cases}
    -\bi\sC(v^{\star}w)\ell(\sM(v^{\star}w)), &\text{if }2\nmid\deg(\sM(v^{\star}w)),\\
    0, &\text{otherwise},
    \end{cases}\quad \text{where } v,w\in\cB_d^{(i,1)};\\
    C^{(i)}_{v,w}&=\begin{cases}
    \sC(v^{\star}w)\ell(\sM(v^{\star}w)), \quad&\text{if }2\nmid\deg(\sM(v^{\star}w)),\\
    0, \quad&\text{otherwise},
    \end{cases}\quad \text{where } v\in\cB_d^{(i,1)},w\in\cB_d^{(i,2)};\\
    D^{(i)}_{v,w}&=\begin{cases}
    -\bi\sC(v^{\star}w)\ell(\sM(v^{\star}w)), &\text{if }2\mid\deg(\sM(v^{\star}w)),\\
    0, &\text{otherwise},
    \end{cases}\quad \text{where } v\in\cB_d^{(i,1)},w\in\cB_d^{(i,2)};\\
    E^{(i)}_{v,w}&=\begin{cases}
    \sC(v^{\star}w)\ell(\sM(v^{\star}w)), \quad&\text{if }2\mid\deg(\sM(v^{\star}w)),\\
    0, \quad&\text{otherwise},\end{cases}\quad \text{where } v,w\in\cB_d^{(i,2)};
    \\
    F^{(i)}_{v,w}&=\begin{cases}
    -\bi\sC(v^{\star}w)\ell(\sM(v^{\star}w)), &\text{if }2\nmid\deg(\sM(v^{\star}w)),\\
    0, &\text{otherwise},
    \end{cases}\quad \text{where } v,w\in\cB_d^{(i,2)}.
    \end{align}
\end{subequations}
For brevity, we will omit the superscript $(i)$ of $A,B,C,D,E,F$ in the rest of this proof.
Then by construction, we have
\begin{equation}\label{sec4:eq1}
    \MM_{d}^{(i)}(\ell)=\begin{bmatrix}
        A+B\bi&C+D\bi\\
        C^{*}-D^{*}\bi&E+F\bi
    \end{bmatrix}.
\end{equation}
For any feasible solution $\ell$ to \eqref{mom}, we define a linear functional $\tilde{\ell}:[\sP_N]_{2d}\rightarrow\C$ by
\begin{equation}
    \tilde{\ell}(u)\coloneqq\begin{cases}\ell(u),&\text{if }\eta(u)=(1,1,1),\\
    0, &\text{otherwise}.\end{cases}
\end{equation}
We next show that $\tilde{\ell}$ is again a feasible solution to \eqref{mom}.
First note that for any $v,w\in\cB_d^{(i)}$, either $\eta(\sM(v^{\star}w))=(1,1,1)$ if $2\mid\deg(\sM(v^{\star}w))$ or $\eta(\sM(v^{\star}w))=(-1,-1,-1)$ if $2\nmid\deg(\sM(v^{\star}w))$, which gives that
\begin{equation}
\MM_{d}^{(i)}(\tilde{\ell})=\begin{bmatrix}
A&D\bi\\
-D^{*}\bi&E
\end{bmatrix}.
\end{equation}
The construction \eqref{sec:eq2} gives $A^*=A$, $E^*=E$, $B^*=-B$, and $F^*=-F$. Applying Lemma~\ref{lm1}, we conclude that \eqref{sec4:eq1} is equivalent to
\begin{equation}
\begin{bmatrix}
A&C&B&D\\
C^*&E&-D^*&F\\
-B&-D&A&C\\
D^*&-F&C^*&E
\end{bmatrix}\succeq0,
\end{equation}
which implies
\begin{equation}
\begin{bmatrix}
A&D\\
D^{*}&E
\end{bmatrix}\succeq0.
\end{equation}
Thus by Lemma \ref{lm2}, we have
\begin{equation}
\MM_{d}^{(i)}(\tilde{\ell})=\begin{bmatrix}
A&D\bi\\
-D^{*}\bi&E
\end{bmatrix}\succeq0.
\end{equation}
Therefore, $\tilde{\ell}$ is feasible for \eqref{mom}.
Moreover, the monomials involved in the Hamiltonian $H$ are all of even degree and we thus have $\tilde{\ell}(H)=\ell(H)$.
This completes the proof.
\end{proof}

\end{document}
